\documentclass[12pt]{article}
\usepackage[utf8]{inputenc}
\usepackage[T1]{fontenc}
\usepackage{hyperref}
\title{Tariana Language Resources for K-12 Curriculum Development}
\author{Compiled by Claude Sonnet 5 (claude-sonnet-5) for Andr\'as Kornai}
\date{}
\begin{document}
\maketitle

\section{Introduction}
Tariana (ISO 639-3: tae; Glottocode: tari1256) is an endangered Arawakan (Maipurean) language of the Vaup\'es basin, Amazonas, Brazil. It is spoken by roughly 100 people, mostly older adults, and speakers are shifting to Tucano. Search results indicate it is the only Arawak language in that linguistic area. Very little classroom material turned up in search; this report lists only what was found.

\section{Primary Educational Resources}
No Tariana primers, textbooks or school materials were found in search results. Wikipedia states that regular Tariana classes have been offered at a school since 2003, but I could not confirm details or find materials.
\begin{itemize}
\item Tariana language (overview, mentions school classes): \url{https://en.wikipedia.org/wiki/Tariana_language}
\end{itemize}

\section{Linguistic and Reference Materials}
\begin{itemize}
\item Aikhenvald, \emph{A Grammar of Tariana, from Northwest Amazonia} (Cambridge University Press, 2003), preface: \url{https://www.cambridge.org/core/books/grammar-of-tariana-from-northwest-amazonia/preface/5D332A41C99B84F1722398329C0290CE}
\item LINGUIST List review of the grammar: \url{https://linguistlist.org/issues/15/606/}
\item Grammar table of contents (PDF): \url{https://www.aikhenvaldlinguistics.com/pdfs/tariana/language/tariana-grammar-contents.pdf}
\item Aikhenvald's Tariana research page (lists a dictionary PDF, which I did not open): \url{https://aikhenvaldlinguistics.com/research/tariana}
\item ``Le tariana'' (PDF, French): \url{https://www.aikhenvaldlinguistics.com/pdfs/tariana/language/le-tariana.pdf}
\end{itemize}

\section{Religious and Community Resources}
None found in search results.

\section{Digital and Audio Resources}
No audio, video or software resources found in search results.

\section{Curriculum Development Context}
The language is used mainly by older adults, so a curriculum would depend on documentation (the Aikhenvald grammar and texts) plus community elders. Local speakers and the school mentioned by Wikipedia should be contacted to learn what materials exist. Tucano and Portuguese are the dominant languages in the area.
\end{document}